\section{Formalización automática}

\subsection{Dos subproblemas}

\begin{enumerate}
    \item \textbf{Formalización de enunciados}: teorema no formal $\to$ enunciado formal
    \item \textbf{Formalización de demostraciones}: demostración no formal $\to$ demostración formal
\end{enumerate}

\subsection{El problema de evaluación de la formalización de enunciados}

Un mismo teorema puede tener varias formulaciones formales equivalentes (por ejemplo, «infinitos números primos» tiene varias escrituras equivalentes), pero comprobar la equivalencia lógica entre dos enunciados formales \textbf{es indecidible en el caso general}.

\subsection{La brecha de razonamiento en la formalización de demostraciones}

Las demostraciones humanas, aun siendo rigurosas, a menudo tienen partes «dejadas al lector»; el modelo debe llenar por sí mismo estas \textbf{brechas de razonamiento}.

\subsection{Caso de estudio: formalización automática en la geometría euclidiana}

Se elige la geometría euclidiana como dominio restringido, aprovechando las particularidades del dominio para volver resoluble un problema originalmente irresoluble.

\subsection{Resumen de la sección}

El reto central de la demostración de teoremas es cómo explorar de manera eficiente un espacio de acciones infinito. Las intuiciones específicas del dominio pueden estructurar el espacio de acciones, pero la generalización sigue siendo un problema abierto.
